% Submitted to The Annals of Statistics on ...
% Revised on ...
\documentclass[a4paper,12pt]{article}
%\documentclass[a4paper,titlepage]{article}

\usepackage{amsmath}
\usepackage{amssymb}
%\usepackage[dvips]{graphicx}
\usepackage{graphicx}
\usepackage[OT4]{fontenc}
%\usepackage[T1]{fontenc}
\usepackage[utf8]{inputenc}
\usepackage{tikz}
\usepackage{url}
\usepackage[numbers]{natbib}
\usepackage{bm}
\emergencystretch=20pt

\usepackage{setspace}
%\doublespacings

\newcommand{\abs}[1]{\left| #1 \right|}
\newcommand{\norm}[1]{\left|\mkern-1mu\left| #1 \right|\mkern-1mu\right|}
\newcommand{\aabs}[1]{\left\| #1 \right\|}
\newcommand{\okra}[1]{\left( #1 \right)}
\newcommand{\kwad}[1]{\left[ #1 \right]}
\newcommand{\klam}[1]{\left\{ #1 \right\}}
\newcommand{\ceil}[1]{\left\lceil #1 \right\rceil}
\newcommand{\floor}[1]{\left\lfloor #1 \right\rfloor}
\newcommand{\dzi}[1]{\left\langle #1 \right\rangle}
\newcommand{\ind}{\mathrel{\perp\mkern-10mu\perp}}
\DeclareMathOperator*{\argmin}{arg\, min}
\DeclareMathOperator*{\argmax}{arg\, max}
\DeclareMathOperator*{\hilberg}{hilb}
\DeclareMathOperator{\essinf}{ess\, inf}
\DeclareMathOperator{\esssup}{ess\, sup}
\DeclareMathOperator{\card}{card}
\DeclareMathOperator{\supp}{supp}
\DeclareMathOperator{\cov}{Cov}
\DeclareMathOperator{\var}{Var}
\DeclareMathOperator{\mean}{\mathbf{E}}
\DeclareMathOperator{\median}{\mathbf{M}}
\newcommand{\peq}{\stackrel{+}{=}}
\newcommand{\pge}{\stackrel{+}{>}}
\newcommand{\ple}{\stackrel{+}{<}}
\newcommand{\boole}[1]{{\bf 1}{\klam{#1}}}

\newtheorem{conjecture}{Conjecture}
\newtheorem{definition}{Definition}
\newtheorem{example}{Example}
\newtheorem{corollary}{Corollary}
\newtheorem{proposition}{Proposition}
\newtheorem{theorem}{Theorem}
\newtheorem{lemma}{Lemma}
\newenvironment*{proof}{\begin{trivlist}\item[]
\noindent\textbf{Proof:}}{$\Box$\par\end{trivlist}}
\newenvironment*{proof*}[1]{\begin{trivlist}\item[]
\noindent\textbf{Proof of #1:}}{$\Box$\par\end{trivlist}}
%\renewcommand{\theenumi}{\roman{enumi}}
%\renewcommand{\labelenumi}{(\theenumi)}
%\renewcommand{\thesection}{\Roman{section}}
%\renewcommand{\thesubsection}{\Roman{section}-\Alph{subsection}}


\newcommand{\bibnamefont}[1]{\MakeUppercase{#1}}


\author{{\L}ukasz D\k{e}bowski\thanks{
    {\L}. D\k{e}bowski is with
    the Institute of Computer Science, Polish Academy of Sciences, 
    ul. Jana Kazimierza 5, 01-248 Warszawa, Poland 
    (e-mail: ldebowsk@ipipan.waw.pl).
  }
  % \thanks{Manuscript received ... ; revised ... .}
  %\thanks{}  
}


\title{Recurrence and repetition times in the case of \\ a stretched
  exponential growth} \date{}

\begin{document}

%\pagestyle{empty}   
\begin{titlepage}
%\titlepage
\maketitle

\begin{abstract}
  By an analogy to the duality between the recurrence time and the
  longest match length, we introduce a quantity dual to the maximal
  repetition length, which we call the repetition time.  Using the
  generalized Kac lemma for successive recurrence times by Chen Moy,
  we sandwich the repetition time in terms of conditional and
  effective min-entropies.
  % The sole assumption is stationarity
  % and ergodicity. The proof is surprisingly short and the claim is
  % fully general in contrast to earlier approaches by Szpankowski and
  % by D\k{e}bowski.
  We discuss the analogy of this result with the
  Wyner-Ziv/Ornstein-Weiss theorem, which sandwiches the recurrence
  time in terms of Shannon entropies. We formulate the respective
  sandwich bounds in a way that applies also to the case of a
  stretched exponential growth observed empirically for natural
  language.
% \end{abstract}
%
% \begin{keywords}
  \\[1em]
  \textbf{Keywords:} stationary processes, recurrence time, maximal
  repetition, R\'enyi entropies
  % \end{keywords}
  \\[1em]
  \textbf{MSC2020:} 60G10, 94A17
\end{abstract}

% %MSC 2020: 
% % 94A29 - Source coding
% % 60G10 - Stationary processes
% % 94A17 - Measures of information, entropy


\end{titlepage}
% \pagestyle{plain}   

% \tableofcontents

\section{Introduction}
\label{secIntroduction}

Using appropriate concepts is often a gate to the simple theory. Some
instance of this rule will be exhibited in this short note. The note
concerns a new stochastic bound for the maximal repetition length and
its newly introduced dual---called the repetition time.  Our
development will be analogous to the well-known duality between the
recurrence time and the longest match length.  Whereas the repetition
time is defined in this paper, there is a long history of
investigation of the recurrence time
\citep{Kac47,WynerZiv89,OrnsteinWeiss93,Kontoyiannis98,Ko12}, the
longest match length \citep{WynerZiv89,OrnsteinWeiss93,Szpankowski93,
  KontoyiannisOthers98,GaoKontoyiannisBienenstock08}, and the maximal
repetition length
\citep{DeLuca99,KolpakovKucherov99a,KolpakovKucherov99,CrochemoreIlie08,
  ErdosRenyi70,ArratiaWaterman89,Shields92b,Shields97,
  Szpankowski93,Szpankowski93a,KontoyiannisSuhov94,
  Debowski11b,Debowski17,Debowski18b}. These concepts appear in
various contexts in ergodic theory, information theory, theoretical
computer science, and quantitative linguistics, as we have indicated
by the references and we will detail further.  Establishing the
asymptotic behavior of these random variables for stationary ergodic
processes is of a theoretical importance since when we are confronted
with empirical data, such results can indicate what kind of a process
may generate these observations.

To be concrete, let $(X_i)_{i\in\mathbb{Z}}$ be an arbitrary
stochastic process over a countable alphabet. The well-researched
recurrence time $R^{(1)}_k$ is the first position in the future on
which a copy of string $X_1^k$ reappears,
\begin{align}
  R^{(1)}_k:=\inf\klam{\infty}\cup\klam{i\ge 1: X_{i+1}^{i+k}=X_1^k}.
\end{align}
By contrast, we define here the repetition time $R^{(2)}_k$ as the
first position in the future on which a copy of any previously seen
string $X_{j+1}^{j+k}$ occurs,
\begin{align}
  R^{(2)}_k:=\inf\klam{\infty}\cup\klam{i\ge 1:
  X_{i+1}^{i+k}=X_{j+1}^{j+k}\text{ for some } 0\le j<i}.
\end{align}
If $(X_i)_{i\in\mathbb{Z}}$ is a process over a $D$-ary alphabet then
$R^{(2)}_k\le D^k$. In any case, $R^{(2)}_k\le R^{(1)}_k$. Precisely,
we have
\begin{align}
  \label{Minimum}
  R^{(2)}_k=\inf_{i\ge 0}\okra{i+R^{(1)}_k\circ T^i}
\end{align}
for the shift automorphism $X_i\circ T=X_{i+1}$, since observing a
repetition is equivalent to observing a recurrence of some particular
string.
% Relationship (\ref{Minimum}) is a simple fact that will have weighty
% consequences in this note.  The motivation for the novel term
% ``impossibility time'' is a rule of thumb that if some rare event
% happens twice then the phenomenon cannot be random.

In parallel, we can discuss variables that are dual to random
variables $R^{(\gamma)}_k$. The longest match length $L^{(1)}_n$ is
\begin{align}
  L^{(1)}_n:=\max\klam{0}\cup\klam{k\ge 1: X_{i+1}^{i+k}=X_1^k
  \text{ for some } 0< i\le n-k},
\end{align}
whereas the maximal repetition length $L^{(2)}_n$ is
\begin{align}
  L^{(2)}_n:=\max\klam{0}\cup\klam{k\ge 1: X_{i+1}^{i+k}=X_{j+1}^{j+k}
  \text{ for some } 0\le j< i\le n-k}.
\end{align}
Both lengths $L^{(\gamma)}_n$ are functions of prefix $X_1^n$ and
$n\ge L^{(2)}_n\ge L^{(1)}_n$. For $\gamma\in\klam{1,2}$, we have the
dualities
\begin{align}
  R^{(\gamma)}_k&=\inf\klam{\infty}\cup\klam{n: L^{(\gamma)}_{n+k}\ge k},
  \\
    L^{(\gamma)}_n&=\max\klam{0}\cup\klam{k: R^{(\gamma)}_k\le n-k}.
\end{align}
Hence $R^{(\gamma)}_k>n\iff L^{(\gamma)}_{n+k}<k$
\citep{WynerZiv89}. Thus the sequence of the dual random variables
carries the same information as the sequence of the primal. Both
$R^{(\gamma)}_k$ and $L^{(\gamma)}_n$ are non-decreasing functions of
$k$ and $n$, respectively.

To comment on the respective theoretical results in the stochastic
setting, let us begin with the recurrence time and the longest match
length.  Let $\log x$ denote the natural logarithm of $x$.  Consider a
discrete random variable $X$. The Shannon entropy is
$H_1(X):=\mean\kwad{-\log P(X)}$. Let $(X_i)_{i\in\mathbb{Z}}$ be a
stationary ergodic process over a finite alphabet. The celebrated
Wyner-Ziv/Ornstein-Weiss (WZ/OW) theorem
\citep{WynerZiv89,OrnsteinWeiss93} asserts
\begin{align}
  \label{WZOW}
  \lim_{k\to\infty} \frac{\log n}{L^{(1)}_n}
  =\lim_{k\to\infty} \frac{1}{k}\log R^{(1)}_k
  =h_1:=\lim_{k\to\infty} \frac{H_1(X_1^k)}{k}
  \text{ a.s.}
\end{align}
Thus the rates of the recurrence time and the longest match length are
essentially dictated by the rate of the Shannon entropy.  The Shannon
entropy rate $h_1$ exists for any stationary process.

Let us proceed to the maximal repetition length and the repetition
time.  For memoryless sources, the maximal repetition length
$L^{(2)}_n$ grows at a similar rate as the longest match length
$L^{(1)}_n$ \citep{ErdosRenyi70,ArratiaWaterman89}. Probably for this
reason, the problem that $L^{(2)}_n$ can grow much faster than
$L^{(1)}_n$ for other stationary processes was first recognized as
late as by \citet{Shields92b,Shields97}. Define the collision entropy
as $H_2(X):=-\log\mean P(X)$, where $H_2(X)\le H_1(X)$ by the Jensen
inequality.  A technically complicated result by
\citet{Szpankowski93,Szpankowski93a}, involving suffix trees, can be
rewritten as
\begin{align}
  \label{Szpankowski}
  \lim_{k\to\infty} \frac{\log n}{L^{(2)}_n}
  =\lim_{k\to\infty} \frac{1}{k}\log R^{(2)}_k
  =h_2:=\lim_{k\to\infty} \frac{H_2(X_1^k)}{k}
  \text{ a.s.}
\end{align}
provided we have a uniform bound on the pointwise mutual information
\begin{align}
  \label{Kfinite}
  \sup_{n,m\ge 1}\esssup
     \frac{P(X_1^{n+m})}{P(X_1^n)P(X_{n+1}^{n+m})}
     < \infty.
\end{align}
Quantity $h_2$ is called the collision entropy rate. Whereas limit
$h_1$ exists for any stationary process, limit $h_2$ exists if
(\ref{Kfinite}) holds additionally. The collision entropy rate can be
substantially smaller than the Shannon entropy rate.

Condition (\ref{Kfinite}) is satisfied by finite-state processes
\citep[Theorem 11.3]{Debowski21} but it does not cover all important
cases---such as stationary sources that do not satisfy the finite
energy condition \citep{Shields92b,Shields97}, do not meet the
Doeblin condition \citep[Theorem 11.7]{Debowski21}, or are not diluted
by random noise \citep{Shields97}, \citep[Theorem 11.6]{Debowski21}.
We became interested in characterizing the maximal repetition length
because of our long-standing interests in statistical language
modeling \citep{Debowski21}.  In fact, for natural language,
experiments suggest that the Shannon entropy rate is positive,
$h_1>0$, \citep{Shannon51,CoverKing78,BehrOthers03} but the mutual
information probably grows like a power-law, called Hilberg's law
\citep{Hilberg90, TakahiraOthers16, KaplanOther20, Debowski21}, which
contradicts condition (\ref{Kfinite}). Moreover, as discovered
empirically in \citep{Debowski12b,Debowski15f} for novels in English,
German, and French, we have a power-law logarithmic growth of the
maximal repetition length
\begin{align}
  L^{(2)}_n\propto \log^\alpha n,\quad \alpha\approx 3.
\end{align}
This is equivalent to the stretched exponential growth of the
repetition time
\begin{align}
  \log R^{(2)}_k\propto k^\beta,\quad \beta\approx 1/3.
\end{align}

The goal of this paper is to analyze this case of a stretched
exponential growth, rectifying and beautifying our earlier theoretical
accounts and cumbersome statements
\citep{Debowski15f,Debowski18b,Debowski21}. We will do it by using the
repetition time $R^{(2)}_k$ rather than the maximal repetition length
$L^{(2)}_n$. The new results are somewhat stronger.

% The new results are simple and general but they have escaped our
% awareness for quite a long time. The reason was that we got stuck to
% bounding $L^{(2)}_n$ rather than $R^{(2)}_k$, gradually losing hope
% in an elegant solution.

A few explanations regarding our framework are in place.  First, the
Hilberg exponent of a sequence $(a_k)_{k\in\mathbb{N}}$ is defined as
\begin{align}
  \label{Exponent}
  \hilberg_{k\rightarrow\infty} a_k
  :=\limsup_{k\rightarrow\infty} \max\klam{0,\frac{\log
  a_k}{\log k}}
  .
\end{align}
It measures the rate of asymptotic power-law growth,
$\hilberg_{k\rightarrow\infty} k^\beta=\max\klam{0,\beta}$. Some
properties of the Hilberg exponent are summarized in \citep[Section
8.1]{Debowski21}. We write succinctly
\begin{align}
  a_k\lesssim b_k\iff
  \hilberg_{k\rightarrow\infty} a_k\le
  \hilberg_{k\rightarrow\infty} b_k,
  \\
  a_k\sim b_k\iff
  \hilberg_{k\rightarrow\infty} a_k=
  \hilberg_{k\rightarrow\infty} b_k.
\end{align}
The Hilberg exponent is useful for characterizing the power-law growth
of mutual information, i.e., Hilberg's law \citep[Chapter
8]{Debowski21}. Hence comes its name.

Second, for a parameter $\gamma\in(0,1)\cup(1,\infty)$, the
R\'enyi-Arimoto conditional entropy of a discrete random variable $X$
given a variable $Y$ is defined as
\begin{align}
  \label{RenyiArimoto}
  H_\gamma(X|Y):=
  -\frac{\gamma}{\gamma-1}\log
  \mean\kwad{\sum_x p(X=x|Y)^{\gamma}}^{1/\gamma}
  ,
\end{align}
whereas for $\gamma\in\klam{0,1,\infty}$, we define
\begin{align}
  H_\gamma(X|Y):=\lim_{\delta\rightarrow\gamma} H_\delta(X|Y).
\end{align}
We also put $H_\gamma(X):=H_\gamma(X|1)$. The special cases are
called:
\begin{itemize}
\item the Hartley entropy
$H_0(X|Y)= \log \esssup\card\klam{x: P(X=x|Y)>0}$,
\item  the Shannon
entropy $H_1(X|Y)=\mean\kwad{-\log P(X|Y)}$, 
\item the collision entropy
$H_2(X|Y)=-2\log \mean\kwad{\sum_{x\in\mathbb{X}} P(X=x|Y)^2}^{1/2}$,
\item the min-entropy $H_\infty(X|Y)=-\log \mean\max_{x} P(X=x|Y)$.
\end{itemize}
We have $H_\gamma(X)\ge H_\delta(X)$ for $\gamma<\delta$, whereas
$H_\gamma(X)\le \frac{\gamma}{\gamma-1}H_\infty(X)$ for $\gamma>1$.
Thus all R\'enyi-Arimoto entropies with $\gamma>1$ are of the same
magnitude.  Moreover, we have a generalized chain rule
\begin{align}
  \label{Chain}
  H_\gamma(X|Y,Z)\le H_\gamma(X|Y)\le H_\gamma(U,X|Y)\le
  H_0(U|Y)+H_\gamma(X|Y,U).
\end{align}
A survey of the R\'enyi-Arimoto entropies can be found in
\citep{FehrBerens14} and also in \citep[Section 9.1]{Debowski21}.

We will also need some modification of the min-entropy.  Let $T$ be a
stationary automorphism of the probability space.  Consider an event
$B$ such that $P(B)>0$. Define the random set
\begin{align}
  \mathcal{C}^B(\omega):=\klam{j\in\mathbb{Z}: T^{-j}\omega\in B}
\end{align}
If $P(B)>0$ then set $\mathcal{C}^B$ contains infinitely many negative
and positive integers almost surely, by the Poincar\'e recurrence
theorem. Then we can arrange the elements of set $\mathcal{C}^B$ into
the unique sorted sequence $(C^B_r)_{r\in\mathbb{Z}}$ defined by
conditions: $\bigcup_{r\in\mathbb{Z}}\klam{C^B_r}=\mathcal{C}$,
$C^B_r>C^B_{r-1}$ for $r\in\mathbb{Z}$, $C^B_1\ge 1$, and
$C^B_0\le 0$. Define the effective min-entropy of a discrete random
variable $X$ as
\begin{align}
  H_+(X):=\log\min_{x} \mean C_2^{(X=x)}.
\end{align}
We will demonstrate this fact.
\begin{theorem}
  \label{theoEffective}
  If automorphism $T$ is ergodic then 
  \begin{align}
    H_+(X)\ge H_\infty(X).
  \end{align}
\end{theorem}

Let us come back to the recurrence and repetition times.  In the body of
the paper, we will prove the following Theorems \ref{theoMainOne} and
\ref{theoMainTwo}, which are analogous but a noticeable novelty can be
attributed only to Theorem \ref{theoMainTwo}.
\begin{theorem}
  \label{theoMainOne}
  Consider a stationary ergodic process $(X_i)_{i\in\mathbb{Z}}$ over
  a finite alphabet. We have a sandwich bound
  \begin{align}
    h_1 k\sim
    \log P(X_1^k|X_{k+1}^\infty)\lesssim
    \log R^{(1)}_k\lesssim
    -\log P(X_1^k)\lesssim
    H_1(X_1^k)
    \text{ a.s.},
  \end{align}
  where the right-most inequality $\lesssim$ becomes the equivalence
  $\sim$ for
  \begin{align}
    \limsup_{k\to\infty} \frac{\sqrt{\var\log P(X_1^k)}}{H_1(X_1^k)}<1.
  \end{align}
\end{theorem}
\emph{Remark:} Theorem \ref{theoMainOne} applies the same sandwich
bound by \citet{Kontoyiannis98} as the simple proof of the WZ/OW
theorem (\ref{WZOW}) presented in \citep[Theorem 9.12]{Debowski21}.
We also recall that $h_1=\mean\kwad{-\log P(X_1|X_2^\infty)}$ for a
stationary process. Consequently if $h_1=0$ then
$P(X_1^k|X_{k+1}^\infty)=1$ almost surely. Quantity
$\var\log P(X_1^k)$ is called varentropy \citep{KontoyiannisVerdu13}.

\begin{theorem}
  \label{theoMainTwo}
  Consider a stationary ergodic process $(X_i)_{i\in\mathbb{Z}}$ over
  a countable alphabet. Define the context length
  \begin{align}
    \label{ContextLength}
    I_k:=\min\klam{i:\log i\ge H_\infty(X_1^k|X_{k+1}^{k+i})}.
  \end{align}
  We have a sandwich bound 
  \begin{align}
    \label{MainTwo}
    H_\infty(X_1^k|X_{k+1}^{k+I_k})\lesssim
    \log I_k\lesssim
    \log R^{(2)}_k\lesssim
    H_+(X_1^k)
    \text{ a.s.}
    ,
  \end{align}
  where the left-most inequality $\lesssim$ becomes the equivalence
  $\sim$ for a finite alphabet.
\end{theorem}
\emph{Remark:} We may substitute
$I_k\to\min\klam{i:\frac{\gamma}{\gamma-1}\log i\ge
  H_\gamma(X_1^k|X_{k+1}^{k+i})}$ for any $\gamma>1$ in the sandwich
bound (\ref{MainTwo}).  In \citep{Debowski18b,Debowski21}, a trivial
bound in terms of Shannon entropy $H_1(X_1^k)$ stemming from
inequality $L^{(2)}_n\ge L^{(1)}_n$ was discussed instead of the
right-most bound in (\ref{MainTwo}). In the previous preprint version
of this article, we claimed erroneously that $\log R^{(2)}_k$ can be
generally upper bounded by $H_\infty(X_1^k)$ rather than
$H_+(X_1^k)$. The phenomenon of burstiness, described in Section
\ref{secRepetition} likely excludes such a possibility.
% By contrast, the upper bound in (\ref{MainTwo}) in terms of the
% effective min-entropy $H_+(X_1^k)$, which is a brand new idea,
% follows by the the generalized Kac lemma for successive recurrence
% times by \citet{ChenMoy59}.
As for the left-most bound in (\ref{MainTwo}), context length $I_k$
can be much shorter than the pessimistic context length $D^k$ for a
$D$-ary alphabet that was proved sufficient in \citep[Theorem
8]{Debowski18b}. Moreover, bounding $R^{(2)}_k$ rather than
$L^{(2)}_n$ leads to more elegant statements than produced in
\citep{Debowski18b,Debowski21}.

This paper intentionally covers theorems rather than examples.
Construction of instances of stationary ergodic processes that satisfy
proportionality $L^{(2)}_n\propto \log^\alpha n$ with exponent
$\alpha>1$ and the Shannon entropy rate $h_1>0$, stated as an open
problem in book \citep{Debowski21}, is relegated to another
article. The essential idea, inspired by Theorem \ref{theoMainTwo}, is
not so difficult---the source should emit sufficiently long words in a
memoryless fashion. However, because of other involved linguistic
concepts, such as Hilberg's law \citep{Hilberg90, TakahiraOthers16,
  KaplanOther20, Debowski21}, and unfinished calculations, we prefer
to publish this detailed result elsewhere.

The organization of this note is as follows. In Section
\ref{secRecurrence}, we recall the known sandwich bound for the
recurrence time. In Section \ref{secRepetition}, we prove a new
sandwich bound for the repetition time. Concluding, in Section
\ref{secFinal}, we specialize these results to the Hilberg exponents.

% Gaps to be filled:

% -- How much can different entropies differ?, Szpankowski's results
% on $L^{(\gamma)}_n$ for $\gamma\ge 3$ suggest that entropy rates
% $h_\gamma$ are equal.


\section{Sandwiching recurrence times}
\label{secRecurrence}

There is a celebrated lemma by \citeauthor{Kac47}. Namely, the
conditional expectation of the recurrence time is the inverse
probability of the initial string.
\begin{lemma}[\citep{Kac47}]
  \label{theoKac}
  For a stationary ergodic process $(X_i)_{i\in\mathbb{Z}}$ over a
  countable alphabet,
  \begin{align}
    \mean\okra{R^{(1)}_k\middle|X_1^k}
    =
    \frac{1}{P(X_1^k)}.
  \end{align}
\end{lemma}
There is also a converse result by \citeauthor{Kontoyiannis98}.
\begin{lemma}[\citep{Kontoyiannis98}]
  \label{theoKontoyiannis}
  For a process $(X_i)_{i\in\mathbb{Z}}$ over a $D$-ary alphabet,
  \begin{align}
    \mean\okra{\frac{1}{R^{(1)}_k
    P(X_1^k|X_2^\infty)}\middle|X_{k+1}^\infty}
    \le
    1+\log D^k.
  \end{align}
\end{lemma}
The origin of this clever but simple inequality comes from an
observation that for a fixed sequence $x_{k+1}^\infty$ and any number
$i\ge 1$, there is at most one string $x_1^k$ such that $R^{(1)}_k=i$
for $X_1^\infty=x_1^\infty$.

The following sandwich bound appears in the simple proof of the WZ/OW
theorem in \citep[Theorem 9.12]{Debowski21}, which was originally
proved using other techniques. We recall our simpler derivation for
the analogy with Section \ref{secRepetition}.
\begin{proposition}
  \label{theoSandwichRecurrence}
  Consider a stationary ergodic process $(X_i)_{i\in\mathbb{Z}}$ over
  a $D$-ary alphabet. Assume that $\sum_{k=1}^\infty \rho_k<\infty$
  for some $\rho_k>0$.  Almost surely, for sufficiently large $k$ we have
\begin{align}
  \label{SandwichRecurrence}
  -\log P(X_1^k|X_{k+1}^\infty)
  +\log \rho_k-\log k
  <
  \log R^{(1)}_k
  <
  -\log P(X_1^k)
  -\log \rho_k
  .
\end{align}
\end{proposition}
\begin{proof}
  As \citet{Kontoyiannis98} has shown, by applying the Markov
  inequality to Lemmas \ref{theoKac} and \ref{theoKontoyiannis}, we
  obtain two complementary bounds
\begin{align}
  \label{ProbRecurrenceI}
  P\okra{R^{(1)}_k\ge \frac{C}{P(X_1^k)}}&\le C^{-1}
  ,
  \\
  \label{ProbRecurrenceII}
  P\okra{R^{(1)}_k\le \frac{C}{P(X_1^k|X_{k+1}^\infty)}}&\le C(1+k\log D)
  .
\end{align}
Setting $C=\rho_k^{-1}$ in (\ref{ProbRecurrenceI}) and $C=k^{-1}\rho_k$ in
(\ref{ProbRecurrenceII}) we obtain (\ref{SandwichRecurrence}) by the
Borel-Cantelli lemma.
\end{proof}

The WZ/OW theorem (\ref{WZOW}) follows by Proposition
\ref{theoSandwichRecurrence} combined with the
Shannon-McMillan-Breiman (SMB) theorem \citep{Shannon48, Breiman57,
  Chung61, Barron85, AlgoetCover88}, which asserts
\begin{align}
  \label{SMB}
  \lim_{k\to\infty}
  \frac{1}{k}\okra{-\log P(X_1^k)}=
  \lim_{k\to\infty}
  \frac{1}{k}\okra{-\log P(X_1^k|X_{k+1}^\infty)}=
  h_1 \text{ a.s.}
\end{align}

\section{Sandwiching repetition times}
\label{secRepetition}

There is a result by \citeauthor{ChenMoy59} that generalizes the Kac
lemma to successive recurrence times. Let $T$ be a stationary
automorphism of the probability space and let
$(C^B_r)_{r\in\mathbb{Z}}$ for $P(B)>0$ be defined as in Section
\ref{secIntroduction}.
% Let $T$ be a stationary
% automorphism of the probability space.  Consider an event $B$ such
% that $P(B)>0$. Define the random set
% \begin{align}
%   \mathcal{C}^B(\omega):=\klam{j\in\mathbb{Z}: T^{-j}\omega\in B}
% \end{align}
% If $P(B)>0$ then set $\mathcal{C}^B$ contains infinitely many negative
% and positive integers almost surely, by the Poincar\'e recurrence
% theorem. Then we can arrange the elements of set $\mathcal{C}^B$ into the
% unique sorted sequence $(C^B_r)_{r\in\mathbb{Z}}$ defined by conditions:
% $\bigcup_{r\in\mathbb{Z}}\klam{C^B_r}=\mathcal{C}$, $C^B_r>C^B_{r-1}$ for
% $r\in\mathbb{Z}$, $C^B_1\ge 1$, and $C^B_0\le 0$.
The successive recurrence times are defined as
\begin{align}
  \label{Successive}
  W^B_r:=C^B_r-C^B_{r-1}.
\end{align}
Observe that $(C^B_0=0)=B$. The result by Chen Moy is as follows.
\begin{lemma}[\citep{ChenMoy59}]
  \label{theoChenMoy}
  Process $(W^B_r)_{r\in\mathbb{N}}$ is stationary with respect to
  conditional measure $P(\cdot|B)$. Moreover, if automorphism $T$
  is ergodic then for $r\in\mathbb{N}$, we have
  \begin{align}
    \mean\okra{W^B_r\middle|B}
    =
    \frac{1}{P(B)}.
  \end{align}
\end{lemma}

Due to the phenomenon of burstiness, this result cannot be
strengthened to the unconditional expectation.
\begin{example}[periodic process with burstiness]
  Consider a periodic process $(X_i)_{i\in\mathbb{Z}}$ in which the
  pattern of $n_1$ ones followed by $n_0$ zeros repeats
  cyclically. Let $B=(X_0=1)$. Then $P(C^B_0=0)=P(B)=n_1/(n_0+n_1)$
  and $P(C^B_0=-i)=1/(n_0+n_1)$ for $i=1,...,n_0$. For $C_0<0$, we
  have $W^B_{kn_1+l}=1+n_0\boole{l=1}$ for
  $k\in\mathbb{N}\cup\klam{0}$ and $l=1,...,n_1$. Hence
\begin{align}
  \mean W^B_{kn_1+l}&=1+\frac{n_0+n_0^2\boole{l=1}}{n_0+n_1}.
\end{align}
For $n_1=1$, we have $\mean W^B_r=1+n_0=1/P(B)$. However, process
$(W^B_r)_{r\in\mathbb{N}}$ is non-stationary for other values of
$n_1$, which correspond to a burst of $1$'s followed a single lull
that is arbitrarily longer than $1/P(X_0=1)$.
\end{example}
We may expect that the phenomenon of burstiness is particularly acute
in statistical models of natural language
\citep{AltmannPierrehumbertMotter09}.

Despite the phenomenon of burstiness, we can provide a simple lower
bound for the expectation of variable $C^B_r$.
\begin{lemma}
  \label{theoUnconditional}
  If automorphism $T$ is ergodic then for $r\in\mathbb{N}$, we have
  \begin{align}
    \mean C^B_r
    \ge
    \frac{r}{2P(B)}.
  \end{align}
\end{lemma}
\begin{proof}
  Since $C^B_m=C^B_0+\sum_{r=1}^m W^B_r$ then by Lemma
  \ref{theoChenMoy}, for $m\in\mathbb{N}$, we obtain
\begin{align}
  \mean\okra{C^B_m\middle|B}
    =
    \frac{m}{P(B)}.
\end{align}
Consider the complement $B'=\Omega\setminus B$ and assume $P(B')>0$.
Specializing Lemma \ref{theoChenMoy}, for $c\in\mathbb{N}$ we obtain
\begin{align}
  \mean\okra{C^{B'}_c\middle|B'}
    =
    \frac{c}{P(B')}.
\end{align}
Comparing sequences $(C^B_m)_{m\in\mathbb{N}}$ and
$(C^{B'}_c)_{c\in\mathbb{N}}$, we notice inequalities $C^B_m-m\ge 0$ and
$C^{B'}_c-c\ge 0$ and equivalence $C^B_m-m<c \iff C^{B'}_c-c>m$. By the
Markov inequality, these yield
\begin{align}
  P(C^B_m-m<c|B')
  &= P(C^{B'}_c-c>m|B')
    \nonumber\\
  &\le \frac{\mean\okra{C^{B'}_c-c\middle|B'}}{m}
    = \frac{c}{m}\okra{\frac{1}{P(B')}-1}.
\end{align}
Hence, we have
\begin{align}
  \mean\okra{C^B_m-m\middle|B'}
  &=
  \int_0^\infty P(C^B_m-m\ge c|B')dc
  \ge
    \int_0^\infty \okra{1-\frac{c}{m}\frac{P(B)}{P(B')}}_+dc
    \nonumber\\
  &=
  m\int_0^{P(B')/P(B)}\okra{1-x\frac{P(B)}{P(B')}}dx
  =
  \frac{m}{2}\frac{P(B')}{P(B)}.
\end{align}
As a result, we derive a lower bound for the unconditional expectation
\begin{align}
  \mean C^B_m
  &=
    P(B)\mean\okra{C^B_m\middle|B}+P(B')\mean\okra{C^B_m\middle|B'}
    \nonumber\\
  &\ge
    m\okra{1+P(B')\okra{1+\frac{1}{2}\frac{P(B')}{P(B)}}}
    =
    \frac{m}{P(B)}\okra{1-\frac{P(B')^2}{2}}
    \nonumber\\
  &\ge
    \frac{m}{2P(B)}.
\end{align}
\end{proof}
In view of Lemma \ref{theoUnconditional} and definition
$H_\infty(X_1^k):=-\log \max_{x_1^k} P(X_1^k=x_1^k)$, we derive
Theorem \ref{theoEffective}.
% inequality
% \begin{align}
%   H_+(X_1^k)\ge H_\infty(X_1^k).
% \end{align}

Let us derive an upper bound for the repetition time.  The following
proposition is a new simple idea.
\begin{proposition}
  \label{theoSandwichRepetitionOne}
  Consider a stationary ergodic process
  $(X_i)_{i\in\mathbb{Z}}$ over a countable alphabet. Assume that
  $\sum_{k=1}^\infty \rho_k<\infty$ for some $\rho_k>0$. Define the
  effective min-entropy
  \begin{align}
    H_+(X_1^k):=\log\min_{x_1^k} \mean C_2^{(X_1^k=x_1^k)}.
  \end{align}
  Almost surely, for sufficiently large $k$ we have
\begin{align}
  \label{SandwichRepetitionOne}
  \log R^{(2)}_k<H_+(X_1^k)-\log\rho_k
  .
\end{align}
\end{proposition}
\emph{Remark:}
% Thus the effective min-entropy is larger than the plain min-entropy.
% We sought for a bound for the maximal repetition length
% intermittently for a few years \citep{Debowski18b,Debowski21}. In
% our perception, this involved lower bounding the probability of a
% union for strongly dependent processes. We tried using the
% Chung-Erd\H{o}s inequality \citep{ChungErdos1952} or other
% second-moment methods like \citet{Szpankowski93,Szpankowski93a} did,
% which led to complicated calculations and assumptions hopeless to
% satisfy. The best that we could achieve was a trivial bound in terms
% of the Shannon entropy rather than the effective min-entropy.
The bound in terms of the effective min-entropy is somewhat difficult
to evaluate. However, we may suppose that if the process does not
exhibit much burstiness and $H_\infty(X_1^k)\sim k^\beta$ then also
$H_+(X_1^k)\sim k^\beta$ and we may have the power-law logarithmic
growth of the maximal repetition length.
\begin{proof}
  Consider event $B=(X_1^k=x_1^k)$ for some string $x_1^k$ such that
  $P(B)>0$ and the corresponding successive recurrence times
  (\ref{Successive}).  We have
  \begin{align}
    R^{(2)}_k\le C_2^{(X_1^k=x_1^k)},
  \end{align}
  since by the time we observe string $x_1^k$ twice, we must have seen
  some repetition of length $k$, certainly.  Fix $x_1^k$ as the most
  probable string of length $k$.  We may bound
  \begin{align}
    P(R^{(2)}_k\ge n)
    &\le \frac{\mean R^{(2)}_k}{n}
      \le \min_{x_1^k} \frac{\mean C^{(X_1^k=x_1^k)}_2}{n}
      = \frac{e^{H_+(X_1^k)}}{n}.
  \end{align}
  Define $n_k:=e^{H_+(X_1^k)}\rho_k^{-1}$. Then we have
  $\sum_{k=1}^\infty P(R^{(2)}_k\ge n_k)<\infty$. Hence inequality
  (\ref{SandwichRepetitionOne}) is obtained by the Borel-Cantelli
  lemma.
\end{proof}

Now we will derive a lower bound for the repetition time.  This
proposition refines \citep[Theorem 9.25]{Debowski21} by using the
weighted conditional entropy rather than the min-entropy conditioned
on the infinite future.
\begin{proposition}
  \label{theoSandwichRepetitionTwo}
  Consider a stationary process $(X_i)_{i\in\mathbb{Z}}$ over a
  countable alphabet. Assume that $\sum_{k=1}^\infty \rho_k<\infty$
  for some $\rho_k>0$. Define the weighted conditional entropy
  \begin{align}
    H_-(X_1^k|X_{k+1}^\infty):=-\log \okra{
    \sum_{i=1}^\infty\frac{e^{-H_\infty(X_1^k|X_{k+1}^{k+i})}}{i(i+1)}}.
  \end{align}
  Almost surely, for sufficiently large $k$ we have
  \begin{align}
    \label{SandwichRepetitionTwo}
    \log R^{(2)}_k> \frac{1}{3}H_-(X_1^k|X_{k+1}^\infty)
    +\frac{1}{3}\log\rho_k
    .
  \end{align}
\end{proposition}
\emph{Remark:} We note that
$H_-(X_1^k|X_{k+1}^\infty)\ge
H_\infty(X_1^k|X_{k+1}^\infty)$. Hence Proposition
\ref{theoSandwichRepetitionTwo} is stronger than \citep[Theorem
9.25]{Debowski21}.
\begin{proof}
  For $i<j$, we obtain 
\begin{align}
  P(X_{i+1}^{i+k}=X_{j+1}^{j+k})
  &=
    \mean
    P(X_{i+1}^{i+k}=X_{j+1}^{j+k}|X_{i+k+1}^{j+k})
    \nonumber\\
  &\le
    \mean \max_{x_1^k} P(X_{i+1}^{i+k}=x_1^k|X_{i+k+1}^{j+k})
    \nonumber\\
  &=
    e^{-H_\infty(X_1^k|X_{k+1}^{k+j-i})}
    .
\end{align}
Observe that $3-2x\le x^{-2}$, so $n-x\le 4n^3x^{-2}/27$.
%$n-x=n/3(3-2(3x/2n))\le n/3(3x/2n)^{-2}=4n^3/27x^2$.
Hence, by the union bound,
\begin{align}
  P(R^{(2)}_k\le n)
  &= P\okra{X_{i+1}^{i+k}=X_{j+1}^{j+k} \text{ for some }
    0\le i<j\le n}
    \nonumber\\
  &\le \sum_{0\le i<j\le n} P(X_{i+1}^{i+k}=X_{j+1}^{j+k})
    \le \sum_{0\le i<j\le n}  e^{-H_\infty(X_1^k|X_{k+1}^{k+j-i})}
    \nonumber\\
  &= \sum_{i=1}^{n-1}
    \frac{n-i}{e^{H_\infty(X_1^k|X_{k+1}^{k+i})}}
    \le \frac{8n^3}{27}
    \sum_{i=1}^\infty\frac{e^{-H_\infty(X_1^k|X_{k+1}^{k+i})}}{i(i+1)}
    \nonumber\\
  &\le \frac{8n^3}{27} e^{-H_-(X_1^k|X_{k+1}^\infty)}.
\end{align}
Put $n_k:=\sqrt[3]{\rho_k e^{H_-(X_1^k|X_{k+1}^\infty)}}$. Then
$\sum_{k=1}^\infty P(R^{(2)}_k\le n_k)<\infty$.  Consequently,
inequality (\ref{SandwichRepetitionTwo}) follows by the
Borel-Cantelli lemma.
\end{proof}

The last proposition in this section is also new although it recycles
some techniques used in universal coding such as \citep[Theorem
6.11]{Debowski21}. The goal is to show that the bound in terms of the
weighted conditional entropy is more precise than the bound proved in
\citep[Theorem 8]{Debowski18b} since the length of conditioning can be
shortened to $I_k$ rather than $D^k$.
\begin{proposition}
  \label{theoSandwichContext}
  Consider a stationary process $(X_i)_{i\in\mathbb{Z}}$ over a
  countable alphabet.  For $I_k$ defined in (\ref{ContextLength}), we
  have
\begin{align}
  \label{SandwichContextOne}
  \log I_k-\log 2&\le H_-(X_1^k|X_{k+1}^\infty)\le
  3\log I_k+\frac{1}{I_k}
  ,
  \\
  \label{SandwichContextThree}
  \log (I_k-1)-H_0(X_i)&\le H_\infty(X_1^k|X_{k+1}^{k+I_k})\le \log I_k
                              .
\end{align}
\end{proposition}
\emph{Remark:} We note that $I_k\le D^k$ and
$H_\infty(X_1^k|X_{k+1}^{k+I_k})\ge H_\infty(X_1^k|X_{k+1}^{k+D^k})$
for a $D$-ary alphabet. Hence Propositions
\ref{theoSandwichRepetitionTwo} and \ref{theoSandwichContext} combined
are stronger than \citep[Theorem 8]{Debowski18b}.
\begin{proof}
First, we derive (\ref{SandwichContextOne}) from
\begin{align}
  H_-(X_1^k|X_{k+1}^\infty)
  &\le
    H_\infty(X_1^k|X_{k+1}^{k+I_k})+\log I_k+\log (I_k+1)\le
    3\log I_k+\frac{1}{I_k},
  \\
  H_-(X_1^k|X_{k+1}^\infty)
  &\ge
    -\log\okra{\okra{1-\frac{1}{I_k}}e^{-H_\infty(X_1^k|X_{k+1}^{k+I_k-1})}
    +\frac{1}{I_k}}
    \nonumber\\
  &\ge -\log\okra{\okra{1-\frac{1}{I_k}}\frac{1}{I_k-1}+\frac{1}{I_k}}
    =\log I_k-\log 2.
\end{align}
Moreover, for a stationary process,
\begin{align}
  0
  &\le
    H_\infty(X_1^k|X_{k+1}^{k+I_k-1})-H_\infty(X_1^k|X_{k+1}^{k+I_k})\le
    H_0(X_i)
\end{align}
since by the chain rule (\ref{Chain}), we have
\begin{align}
  H_\infty(Y,X|Z,V)\le H_\infty(Y,X|Z)
  &\le H_0(X|Z)+H_\infty(Y|X,Z)
    \nonumber\\
  &\le H_0(X)+H_\infty(U,Y|X,Z). 
\end{align}
Hence we obtain (\ref{SandwichContextThree}).
\end{proof}

\section{Final touch}
\label{secFinal}

Let us recall some properties of Hilberg exponents.  The probabilistic
results of \citep[Section 8.1]{Debowski21} can be summarized as
follows.
\begin{definition}
  Let $T$ be an automorphism of the probability space.  A sequence of
  random variables $(J_k)_{k\in\mathbb{N}}$ is called $T$-increasing
  if
  \begin{align}
    J_{k+1}\ge J_k\circ T,J_k\ge 0.
  \end{align}
\end{definition}
Let $\median X:=\sup \klam{r: P(X<r)\le \frac{1}{2}}$ denote the
median of a variable $X$.
\begin{lemma}
  \label{theoHilberg}
  For a stationary ergodic automorphism $T$ and a $T$-increasing
  sequence $(J_k)_{k\in\mathbb{N}}$, we have
  \begin{align}
    \hilberg_{k\rightarrow\infty} \median J_k
    \le \hilberg_{k\rightarrow\infty} J_k 
    \le \hilberg_{k\rightarrow\infty} \mean J_k
    \text{ a.s.}, 
  \end{align}
  where random variable $\hilberg_{k\rightarrow\infty} J_k$ is almost
  surely constant. Moreover, the inequalities become the equalities if
  \begin{align}
    \limsup_{k\rightarrow\infty} \frac{\sqrt{\var J_k}}{\mean J_k}<1.
  \end{align}
\end{lemma}
\begin{proof}
  All claims but the last one follow by \citep[Theorems 8.3--8.5]{Debowski21}.
  The last claim is a slight strengthening of \citep[Theorem
  8.6]{Debowski21} and in a similar fashion it follows by the
  median-mean inequality $\abs{\median X-\mean X}\le \sqrt{\var X}$,
  see e.g.\ \citep[Theorem 3.17]{Debowski21}.
\end{proof}

We also make this simple observation.
\begin{proposition}
  \label{theoIncreasing}
  Consider the shift automorphism $X_i\circ T=X_{i+1}$. Then sequences
  $(L^{(\gamma)}_n)_{n\in\mathbb{N}}$,
  $(R^{(\gamma)}_k)_{k\in\mathbb{N}}$,
  $(-\log P(X_1^k))_{k\in\mathbb{N}}$, and
  $(-\log P(X_1^k|X_{k+1}^\infty))_{k\in\mathbb{N}}$ are
  $T$-increasing.
\end{proposition}
Thus Theorem \ref{theoMainOne} follows by Propositions
\ref{theoSandwichRecurrence} and \ref{theoIncreasing} and Lemma
\ref{theoHilberg}, whereas Theorem \ref{theoMainTwo} follows by
Propositions
\ref{theoSandwichRepetitionOne}--\ref{theoSandwichContext}. Quod erat
demonstrandum.

% \bibliographystyle{IEEEtran}
\bibliographystyle{abbrvnat}
%\bibliographystyle{elsart-harv}

\bibliography{0-journals-abbrv,0-publishers-abbrv,ai,mine,tcs,ql,nlp,books}

\end{document}

%%% Local Variables: 
%%% mode: latex
%%% mode: flyspell
%%% TeX-master: t
%%% TeX-PDF-mode: t
%%% ispell-local-dictionary: "american"
%%% End: 

%  LocalWords:  differentiable


% Suppose that $n_1\ge 2$. Then
% \begin{align}
%   \mean C_2=\mean C_0+\mean W_1+\mean W_2=
%   2+\frac{n_0(n_0+3)}{2(n_0+n_1)}.
% \end{align}

% \begin{align}
%   \mean W_r=\int_0^\infty P(W_r\ge n) dn.
% \end{align}
% Consider events $B_0=B$ and
% $B_n:=T^{-n}B\cap T^{-n-1}B'\cap...\cap B'$ for $n\ge 1$. Class
% $\klam{B_n}_{n\ge 0}$ is a partition. We have $C_m\le C_{n,m}-n$ and
% \begin{align}
%   \mean C_m=\sum_{n\ge 0} P(B_n)\mean\okra{C_m\middle|B_n}\le
%   \sum_{n\ge 0} (1-nP(B_n))=\infty+\mean C_0.
% \end{align}

% Bursts and lulls of words. If there are bursts and lulls then
% $\mean\okra{C_m\middle|B}\le \mean\okra{C_m\middle|B'}$. How to
% quantify bursts and lulls of a given word? Does this measure tend to
% zero for the growing word length in the stationary ergodic case?

% Let $\beta$ be a partition. We have
% \begin{align}
%   \mean \min_{\hat B\in\beta} \hat C_m
%   =
%   \sum_{B\in\beta} P(B)\mean\okra{\min_{\hat B\in\beta}
%   \hat C_m\middle|B}
%   \le
%   \sum_{B\in\beta} P(B)\mean\okra{C_m\middle|B}
%   =
%   m\card \beta.
% \end{align}

% We have
% \begin{align}
%   \lim_{m\to\infty}\frac{1}{m}\sum_{i=1}^m W_i
%   =\okra{\lim_{k\to\infty}\frac{1}{k}\sum_{j=1}^k I_B\circ T^j}^{-1}
%   =\frac{1}{P(B)}
%   \text{ $P$-almost surely}.
% \end{align}
% Then $\mean C_m=\mean C_0+m/P(B)$. Hence 
% \begin{align}
%   0\ge \mean C_0\ge -\frac{1}{2}\okra{\frac{1}{P(B)}-2+P(B)}.
% \end{align}
% Symmetry broken and restored!

% \begin{align}
%   \mean(C_m\wedge C'_c)
%   \le
%   P(B)\mean\okra{C_m\middle|B}+P(B')\mean\okra{C'_c\middle|B'}
%   \le m+c.
% \end{align}
% Put $c=\ceil{m/P(B)}$.
% % \begin{align}
% %   C'_c<C_m\iff C'_c-c<C_m-c\iff
% %   C_{C_m-c}-C_m+c>c\iff C_{C_m-c}>C_m.
% % \end{align}
% \begin{align}
%   C_m+C'_c=C_m\wedge C'_c+C_m\vee C'_c.
% \end{align}
% \begin{align}
%   C_m\ge m, \quad C'_c\ge c, \quad
%   C_m\wedge C'_c\le m+c, \quad C_m\vee C'_c\ge m+c.
% \end{align}

% Some measure of burstiness of event $B$ is the difference
% \begin{align}
%   \sup_{r\in\mathbb{N}}\mean W_r-\frac{1}{P(B)}\ge 0.
% \end{align}
% For $C_0=-i$, we have $W_1\ge 1+i$.
% \begin{align}
%   \lim_{m\to\infty}\frac{1}{m}\sum_{i=1}^m \mean W_i=\frac{1}{P(B)}.
% \end{align}

% Let us try a different way. By symmetry, we have
% $C_m-m>c \iff C'_c-c<m$. Hence
% \begin{align}
%   \mean\okra{C_m-m\middle|B'}\le
%   1+\sum_{c=1}^\infty P\okra{C_m-m>c\middle|B'}=
%   1+\sum_{c=1}^\infty P\okra{C'_c-c<m\middle|B'}.
% \end{align}
% In particular,
% \begin{align}
%   \mean\okra{C_1\middle|B'}
%   \le
%   2+\sum_{c=1}^\infty P\okra{C'_c=c\middle|B'}
%   =
%   2+\sum_{c=1}^\infty P\okra{\bigcap_{j=1}^c T^j B'\middle|B'}
%   .
% \end{align}

% For $\epsilon<\mean X$, we have the Paley-Zygmund inequality
% \begin{align}
%   P(X<\epsilon)\le \frac{\var X}{\var X+(\mean X-\epsilon)_{+}^2}.
% \end{align}

